\chapter{Classes and Objects}\label{ch:classes}

Every Dylan object is a \emph{direct instance} of exactly one class, its
class, and an \emph{indirect instance} of all that class's superclasses; a
\term{general instance} of a class is a direct or an indirect instance.
Classes describe the structure of their instances, the \term{slots} that
hold an object's state, and together with generic functions they determine
the instances' behavior. Classes are themselves objects, instances of
\code{<class>}, and can be passed around and examined like any other.

This chapter describes defining classes and their slots, making and
initializing instances, the functions that examine classes, singletons and
sealing, and the classes built into the language. The standard is
\dylanbook{chapters 2, 6 and 11}. Dispatch on classes is described in
Chapter~\ref{ch:functions}.

\section{Defining classes: define-class}\label{sec:classes-define-class}

\index{define-class@\code{define-class}}
\begin{quote}
\code{(define-class \meta{name} (\meta{superclass} \ldots) \meta{slot-spec} \ldots)}
\end{quote}

\noindent \code{define-class} creates a class and stores it in the module
variable \meta{name}, which is made read-only. The superclass expressions
are evaluated and must give at least one class, with no class twice; the
new class inherits from all of them. Each \meta{slot-spec} adds a slot
(Section~\ref{sec:classes-slots}). At the interactive prompt
\code{define-class} returns the name.

\begin{dylan}
(define-class <point> (<object>)
  (x init-value: 0 init-keyword: x:)
  (y init-value: 0 init-keyword: y:))
(define p (make <point> x: 3))
(x p)                          ; => 3
(y p)                          ; => 0
(set! (y p) 4)                 ; => 4
(instance? p <point>)          ; => #t
\end{dylan}

\noindent By convention class names begin and end with angle brackets. A
class that needs no other superclass names \code{<object>}, the root of
the class hierarchy.

In the Common Lisp implementation, defining a class whose name already
holds a class modifies that class in place, as the 1992 manual says, so
that everything referring to it sees the new definition. In the compiler to
C a class may be defined only once.

\subsection{Multiple inheritance}\label{sec:classes-multiple-inheritance}

A class may have several direct superclasses. It inherits the slots of all
of them, and methods specialized on any of them apply to its instances.
The order of the superclasses matters: it determines the class's
\term{class precedence list}, which orders the class and all its
superclasses from most to least specific, and so decides which method runs
when methods on two unrelated superclasses both apply
(Section~\ref{sec:functions-cpl}).

\begin{dylan}
(define-class <named> (<object>)
  (name init-keyword: name: init-value: "anonymous"))
(define-class <located> (<object>)
  (place init-keyword: place: init-value: 'nowhere))
(define-class <landmark> (<named> <located>))
(define l (make <landmark> name: "Big Ben" place: 'london))
(list (name l) (place l))      ; => ("Big Ben" london)
(all-superclasses <landmark>)
  ; => ({the class <landmark>} {the class <named>}
  ;     {the class <located>} {the class <object>})
\end{dylan}

\noindent A class cannot be its own superclass, directly or indirectly. If
the superclasses' orderings cannot all be satisfied at once, no class
precedence list exists and defining the class is an error.

\section{Slots}\label{sec:classes-slots}

A slot is a place in an object that holds a value. A slot is not read by
name; it is read and written by \term{getter} and \term{setter} generic
functions, to which \code{define-class} adds methods. A slot named
\code{x} gets a getter method on the generic function \code{x} and a setter
method on \code{(setter x)}, and \code{(set! (x p) 10)} calls the setter.
Since accessors are generic functions, a slot looks to the rest of the
program exactly like a computed property, and a subclass can replace a
stored slot with a computed one without changing any caller.

\subsection{Slot specifications}

A slot specification takes one of three forms:

\begin{quote}
\meta{slot-name}\\
\code{(\meta{slot-name} \meta{keyword} \meta{value} \ldots)}\\
\code{(\meta{keyword} \meta{value} \ldots)}
\end{quote}

\noindent The slot name is a symbol and is not evaluated. It supplies the
default getter name, \meta{slot-name}, and the default setter name,
\code{(setter \meta{slot-name})}. In the third form there is no slot name,
and \code{getter:} must be given. Each option may be given at most once.

\begin{description}
\item[\code{getter:} \meta{name}] The variable to which the getter method
  is added. Not evaluated.
\item[\code{setter:} \meta{name}] The variable to which the setter method
  is added. Not evaluated. \code{setter: \#f} gives the slot no setter.
\item[\code{type:} \meta{type}] A class or singleton that the slot's values
  must be instances of; storing anything else is an error. Evaluated once.
  The default is \code{<object>}.
\item[\code{init-value:} \meta{value}] The slot's initial value, used for
  every new instance. Evaluated once, when the class is defined.
\item[\code{init-function:} \meta{function}] A function of no arguments,
  called to compute the initial value for each new instance. Evaluated once,
  to give the function. Cannot be combined with \code{init-value:}.
\item[\code{init-keyword:} \meta{keyword}] A keyword with which
  \code{make} may be given the initial value. Not evaluated.
\item[\code{required-init-keyword:} \meta{keyword}] Like
  \code{init-keyword:}, but \code{make} must be given it. Cannot be
  combined with \code{init-keyword:}, \code{init-value:} or
  \code{init-function:}.
\item[\code{allocation:} \meta{kind}] Where the slot's value is kept:
  \code{instance}, \code{class}, \code{each-subclass}, \code{constant} or
  \code{virtual} (Section~\ref{sec:classes-allocation}). Not evaluated. The
  default is \code{instance}.
\end{description}

\begin{dylan}
(define-class <rectangle> (<object>)
  (top type: <integer> init-value: 0 init-keyword: top:)
  (left type: <integer> init-value: 0 init-keyword: left:)
  (bottom type: <integer> init-value: 100 init-keyword: bottom:)
  (right type: <integer> init-value: 100 init-keyword: right:))
(define r (make <rectangle> top: 50 left: 50))
(top r)                        ; => 50
(bottom r)                     ; => 100
(set! (bottom r) 55)           ; => 55
(set! (bottom r) 'foo)         ; => error: foo is not an instance of <integer>
\end{dylan}

\noindent The getter and setter options let a slot have accessor names that
differ from its slot name, or no setter at all:

\begin{dylan}
(define-class <account> (<object>)
  (number required-init-keyword: number: setter: #f)
  (balance init-value: 0
           getter: current-balance setter: (setter current-balance))
  (getter: owner-name init-keyword: owner:))
(define a (make <account> number: 1001 owner: "Ann"))
(number a)                     ; => 1001
(current-balance a)            ; => 0
(owner-name a)                 ; => "Ann"
(make <account>)               ; => error: number: is required
\end{dylan}

\subsection{Initial values}\label{sec:classes-slot-init}

When \code{make} creates an instance, each slot is initialized from the
first of these that applies: the init-keyword's argument to \code{make}, if
given; the init-value; the result of calling the init-function. The
init-function is not called when the keyword is given. A slot with none of
these is left \term{uninitialized}, and reading it is an error until it has
been set. The initial value is stored directly into the slot; the setter
generic function is not called.

An init-value is evaluated once, when the class is defined, and the same
object becomes the initial value of every instance. To give each instance
its own fresh object, such as an empty collection, use an init-function:

\begin{dylan}
(define-class <bag> (<object>)
  (shared init-value: (make <stretchy-vector>))
  (own init-function: (method () (make <stretchy-vector>))))
(define b1 (make <bag>))
(define b2 (make <bag>))
(id? (shared b1) (shared b2))  ; => #t
(id? (own b1) (own b2))        ; => #f
\end{dylan}

\index{slot-initialized?@\code{slot-initialized?}}
\code{(slot-initialized? \meta{instance} \meta{getter})} tells whether the
slot that \meta{getter} reads in \meta{instance} has a value. It is an error
if \meta{getter} reads no slot of the instance. There is no way to make a
slot uninitialized again.

\begin{dylan}
(define-class <memo> (<object>) (note))
(define m (make <memo>))
(slot-initialized? m note)     ; => #f
(note m)                       ; => error: the slot note is not initialized
(set! (note m) "remember")     ; => "remember"
(slot-initialized? m note)     ; => #t
\end{dylan}

\subsection{Allocation}\label{sec:classes-allocation}

\begin{description}
\item[\code{instance}] Each instance has its own value. This is the
  default.
\item[\code{class}] There is one value, shared by all the general instances
  of the class, its subclasses' instances included. Setting it through one
  instance changes it for all.
\item[\code{each-subclass}] The class has one value, shared by its direct
  instances, and each subclass has a value of its own, shared by the
  subclass's direct instances.
\item[\code{constant}] The value is the init-value, which must be given,
  and it cannot be changed: the slot has no setter.
\item[\code{virtual}] No storage is allocated. The program defines methods
  on the getter, and on the setter if it wants one, that compute the value;
  \code{define-class} only makes sure the generic functions exist.
\end{description}

\noindent A \code{class} or \code{each-subclass} slot can be initialized
only with \code{init-value:}; without one it starts uninitialized, like
any other slot, until it is set; a \code{virtual} slot takes no initialization
options at all, and \code{slot-initialized?} is always true of it.

\begin{dylan}
(define-class <vehicle> (<object>)
  (count allocation: class init-value: 0)
  (kind allocation: each-subclass init-value: 'vehicle)
  (wheels allocation: constant init-value: 4)
  (description allocation: virtual))
(define-class <truck> (<vehicle>))
(define-method description ((v <vehicle>))
  (list (kind v) (wheels v)))
(define bike (make <vehicle>))
(define truck (make <truck>))
(set! (count bike) 2)          ; => 2
(count truck)                  ; => 2
(set! (kind truck) 'truck)     ; => truck
(kind bike)                    ; => vehicle
(kind (make <truck>))          ; => truck
(description truck)            ; => (truck 4)
(slot-initialized? bike description)   ; => #t
\end{dylan}

\noindent A virtual slot can also give an inherited stored slot a new
behavior in a subclass, a \term{filtered slot}: the subclass declares the
slot virtual for documentation and defines getter and setter methods that
call \code{next-method} to reach the stored value.

\begin{dylan}
(define-class <view> (<object>)
  (position init-value: 0 init-keyword: position:))
(define-class <displaced-view> (<view>)
  (offset init-value: 10))
(define-method position ((v <displaced-view>))
  (+ (next-method) (offset v)))
(define-method (setter position) ((v <displaced-view>) new)
  (next-method v (- new (offset v))))
(define dv (make <displaced-view> position: 5))
(position dv)                  ; => 15
(set! (position dv) 100)       ; => 100
(position dv)                  ; => 100
\end{dylan}

\subsection{Inherited slots}\label{sec:classes-inherited-slots}

Two slots are the same slot if they have the same getter. A subclass may
respecify a slot it inherits, to give it, for example, a different initial
value. The 1992 manual asks that the new specification repeat every feature
the inherited one specified (its type, an initial value, its init-keyword,
a non-instance allocation and a setter), with a type no wider than the
inherited one, and both implementations report a respecification that
leaves one out (``redefining an inherited slot must respecify its
init-keyword:''). The exception is a slot redeclared \code{virtual}, a
filtered slot (below), which respecifies nothing: the inherited slot keeps
its storage, its type and its initialization, so \code{make} still takes
its init-keyword. A class that inherits the same slot from two
superclasses with different specifications must respecify it.

\begin{dylan}
(define-class <base> (<object>)
  (color init-value: 'red init-keyword: color:))
(define-class <derived> (<base>)
  (color init-value: 'blue init-keyword: color:))
(color (make <base>))                 ; => red
(color (make <derived>))              ; => blue
(color (make <derived> color: 'green)) ; => green
\end{dylan}

\section{Making instances: make and initialize}\label{sec:classes-make}

\index{make@\code{make}}\index{initialize@\code{initialize}}
\begin{quote}
\code{(make \meta{class} \meta{keyword} \meta{value} \ldots)}
\end{quote}

\noindent \code{make} returns a general instance of \meta{class}. For a
class the program defines, it allocates a new instance, initializes its
slots (Section~\ref{sec:classes-slot-init}), then calls the generic function
\code{initialize} with the instance and all the keyword arguments, and
returns the instance; what \code{initialize} returns is ignored. Making an
instance of an abstract class, such as \code{<number>} or
\code{<collection>}, is an error. Built-in instantiable classes have
\code{make} methods of their own, described with them
(Chapter~\ref{ch:collections}).

The keywords \code{make} accepts are the init-keywords of the class's slots,
inherited slots included, and the keywords accepted by the applicable
\code{initialize} methods. An \code{initialize} method with \code{\#rest}
accepts any keyword. Any other keyword is an error, as the manual
specifies: the compiler to C reports \code{keyword argument not accepted:
y:}, the Common Lisp implementation \code{make -{}- illegal keyword
argument y:}.

Programs customize instance creation by adding methods to
\code{initialize}, not to \code{make}. An \code{initialize} method normally
calls \code{next-method} so that the superclasses' methods run too; the
method on \code{<object>} does nothing. When \code{initialize} runs, the
slots already have their initial values.

\begin{dylan}
(define-class <shape> (<object>)
  (name init-keyword: name: init-value: "shape")
  (area))
(define-class <square> (<shape>)
  (side init-keyword: side: init-value: 1))
(define-method initialize ((s <shape>) #rest options)
  (next-method)
  (print (list 'initializing (name s))))
(define-method initialize ((s <square>) #rest options)
  (next-method)
  (set! (area s) (* (side s) (side s))))
(define sq (make <square> side: 3 name: "sq"))
                               ; prints: (initializing "sq")
(area sq)                      ; => 9
\end{dylan}

\noindent The \code{<square>} method above takes \code{\#rest} rather than
\code{\#key side}: a parameter named \code{side} would shadow the getter
\code{side} within the method. To receive a keyword under another name,
write \code{\#key (side: the-side)} (Section~\ref{sec:functions-keywords}).

\section{Adding slots: define-slot}\label{sec:classes-define-slot}

\index{define-slot@\code{define-slot}}
\begin{quote}
\code{(define-slot \meta{owner} [\meta{slot-name}] \meta{keyword} \meta{value} \ldots)}
\end{quote}

\noindent \code{define-slot} adds a slot to an existing class or to a
single object. \meta{owner} is evaluated and must be a class or a
singleton; the slot name and options are those of a slot specification in
\code{define-class}. If the owner is a class, all its instances get the
slot. If it is a singleton, only that one object does, and its allocation
must be \code{instance}, \code{constant} or \code{virtual}, without
\code{init-keyword:} or \code{required-init-keyword:}.

\begin{dylan}
(define-class <pt> (<object>) (px init-value: 0))
(define p (make <pt>))
(define-slot (singleton p) note init-value: "only me")
(note p)                       ; => "only me"
(note (make <pt>))             ; => error: no applicable method
(define x (list 1 2))
(define-slot (singleton x) tag init-value: 'special)
(set! (tag x) 'changed)        ; => changed
(tag x)                        ; => changed
(define-slot (singleton 'cup) volume allocation: constant init-value: 8)
(volume 'cup)                  ; => 8
(define-slot <pt> weight init-value: 10)
(weight p)                     ; => 10
\end{dylan}

\noindent As the examples show, the object may be of any class, built-in or
not, and the slot's accessors are generic functions with a method
specialized on the singleton.

In the Common Lisp implementation \code{define-slot} takes effect when it is
evaluated, and instances that already exist gain the slot. In the compiler
to C the slot is part of the program from the start: a slot defined on a
class is in every instance ever made, and is listed by
\code{slot-descriptors} even before the \code{define-slot} form is reached.
The compiler to C also requires a class owner to be a class the program
itself defines, known when compiling, rather than a built-in class or a
value computed at run time; a singleton owner may be any object.

\section{Examining classes and objects}\label{sec:classes-reflection}

\begin{description}
\item[\code{(object-class \meta{object})}]
  \index{object-class@\code{object-class}}
  returns the class of which \meta{object} is a direct instance.
\item[\code{(instance? \meta{object} \meta{type})}]
  \index{instance?@\code{instance?}}
  returns true if \meta{object} is a general instance of the class
  \meta{type}, or is the object of the singleton \meta{type}.
\item[\code{(subclass? \meta{class1} \meta{class2})}]
  \index{subclass?@\code{subclass?}}
  returns true if \meta{class1} is \meta{class2} or a direct or indirect
  subclass of it.
\item[\code{(all-superclasses \meta{class})}]
  \index{all-superclasses@\code{all-superclasses}}
  returns the class precedence list: \meta{class} first, \code{<object>}
  last.
\item[\code{(direct-superclasses \meta{class})}]
  returns the superclasses named in the class's definition, in order.
\item[\code{(direct-subclasses \meta{class})}]
  returns the classes that name \meta{class} as a direct superclass.
\end{description}

\begin{dylan}
(define-class <animal> (<object>))
(define-class <dog> (<animal>))
(define rex (make <dog>))
(object-class rex)             ; => {the class <dog>}
(instance? rex <animal>)       ; => #t
(instance? 3 <real>)           ; => #t
(instance? 3 (singleton 3))    ; => #t
(subclass? <dog> <animal>)     ; => #t
(subclass? <dog> <dog>)        ; => #t
(direct-superclasses <dog>)    ; => ({the class <animal>})
(direct-subclasses <animal>)   ; => ({the class <dog>})
(all-superclasses <dog>)
  ; => ({the class <dog>} {the class <animal>} {the class <object>})
\end{dylan}

\noindent None of the sequences these functions return should be modified.
The two implementations give different answers for some built-in objects.
In the compiler to C, \code{\#t} and \code{\#f} are instances of an
internal class that prints as \code{\{the class <boolean>\}} but is not the
value of any variable a program can name, while in the Common Lisp
implementation their class is \code{<object>}; and the Common Lisp implementation seals the
built-in classes that cannot be subclassed, such as \code{<integer>} and
\code{<pair>}, so \code{all-superclasses} and \code{direct-superclasses} of
them return empty sequences (Section~\ref{sec:classes-sealing}).

\subsection{Slot descriptors}\label{sec:classes-slot-descriptors}

A \term{slot descriptor}, of class \code{<slot-descriptor>}, describes one
slot. Programs seldom need them, but debuggers and inspectors do. Both
implementations provide the slot-descriptor functions of
\dylanbook{chapter 11}. (Design note 9 later removed them from the
language.)

\begin{description}
\item[\code{(slot-descriptors \meta{class})}] all the slots of the class,
  inherited ones included.
\item[\code{(slot-descriptor \meta{instance} \meta{getter})}] the slot that
  \meta{getter} reads in \meta{instance}, or \code{\#f}.
\item[\code{(slot-getter \meta{sd})}, \code{(slot-setter \meta{sd})}] the
  getter and setter generic functions; \code{slot-setter} returns
  \code{\#f} for a slot without a setter.
\item[\code{(slot-type \meta{sd})}] the type, \code{<object>} if none was
  given.
\item[\code{(slot-allocation \meta{sd})}] the allocation, a symbol.
\item[\code{(init-value \meta{sd})}] two values: the init-value and
  \code{\#t}, or \code{\#f} and \code{\#f} if there is none.
\item[\code{(init-function \meta{sd})}, \code{(init-keyword \meta{sd})}]
  the init-function or init-keyword, or \code{\#f}.
\item[\code{(slot-value \meta{instance} \meta{sd})}] the value of the slot
  in the instance. Its setter, used as
  \code{(set! (slot-value \meta{instance} \meta{sd}) \meta{value})},
  sets it, checking the slot's type. Neither works on a virtual slot.
\end{description}

\begin{dylan}
(define-class <cell> (<object>)
  (contents type: <integer> init-value: 1 init-keyword: contents:))
(define c (make <cell> contents: 3))
(size (slot-descriptors <cell>))       ; => 1
(define sd (slot-descriptor c contents))
(id? (slot-getter sd) contents)        ; => #t
(id? (slot-setter sd) (setter contents)) ; => #t
(slot-type sd)                         ; => {the class <integer>}
(slot-allocation sd)                   ; => instance
(init-value sd)                        ; => 1, #t
(init-keyword sd)                      ; => contents:
(init-function sd)                     ; => #f
(slot-value c sd)                      ; => 3
(set! (slot-value c sd) 4)             ; => 4
(contents c)                           ; => 4
(slot-descriptor c object-class)       ; => #f
\end{dylan}

\subsection{Copying and converting}

\code{(shallow-copy \meta{object})} returns a new object with the same
contents, made with \code{make} on \code{(class-for-copy \meta{object})},
which by default is the object's class. The library has methods for
collections. For other objects the 1992 manual leaves the method to the
programmer, and that is so in the Common Lisp implementation; the compiler
to C has a default method that copies any instance's slots.

\code{(as \meta{class} \meta{object})} converts \meta{object} to
\meta{class}, returning it unchanged if it is already an instance. Its
methods for numbers, characters, symbols and collections are described in
Chapters~\ref{ch:numbers} and~\ref{ch:collections}.

\section{Singletons}\label{sec:classes-singletons}

A singleton is a specializer for one object (Section~\ref{sec:functions-singletons}).
\code{(singleton \meta{object})} and \code{(make <singleton> object:
\meta{object})} both make one. With methods on singletons a program can give
an individual object behavior of its own, and with \code{define-slot}
(Section~\ref{sec:classes-define-slot}) state of its own, without defining a
class for it. A method on a singleton is more specific than any method on
the object's classes.

%! lisp-only
\begin{dylan}
(define-class <window> (<object>))
(define main-window (make <window>))
(define-method title ((w <window>)) "untitled")
(define-method title ((w (singleton main-window))) "Main")
(title (make <window>))        ; => "untitled"
(title main-window)            ; => "Main"
\end{dylan}

\noindent The compiler to C does not accept that program: there, a
singleton used as a method's specializer must be of a literal, a class, or
a variable whose definition gives it a literal value, because the
specializer must be known when compiling
(Section~\ref{sec:functions-singletons}). A singleton of an instance made at
run time can still own slots, through \code{define-slot}, and singletons of
symbols, keywords, numbers and characters work everywhere:

\begin{dylan}
(define-method plural ((word <symbol>))
  (as <symbol> (concatenate (as <string> word) "s")))
(define-method plural ((word (singleton 'mouse))) 'mice)
(plural 'cat)                  ; => cats
(plural 'mouse)                ; => mice
(make <singleton> object: 4)   ; => {singleton 4}
\end{dylan}

\section{Sealing}\label{sec:classes-sealing}

\index{seal@\code{seal}}\index{make-read-only@\code{make-read-only}}
\code{(seal \meta{class})} seals a class: no subclass of it can be created
any more, and \code{all-superclasses}, \code{direct-superclasses} and
\code{slot-descriptors} return empty sequences for it. Sealing changes
nothing about existing instances or methods. \code{(make-read-only
\meta{class})} forbids changing the class's slots and superclasses without
forbidding subclasses. Both return the class.

\begin{dylan}
(define-class <a> (<object>) (a-slot init-value: 1))
(define-class <b> (<a>))
(seal <b>)                     ; => {the class <b>}
(all-superclasses <b>)         ; => ()
(direct-superclasses <b>)      ; => ()
(a-slot (make <b>))            ; => 1
(instance? (make <b>) <a>)     ; => #t
\end{dylan}

\noindent In the Common Lisp implementation, defining a subclass of a sealed
class is then an error. In the compiler to C all of a program's classes are
created when it is compiled, before \code{seal} runs, so \code{seal} affects
only what the reflective functions return afterwards; and there
\code{make-read-only} on a class does exactly what \code{seal} does. A
built-in class that cannot be subclassed, such as \code{<integer>} or
\code{<list>}, is reported when compiling, if a program tries.

\section{Changing classes at run time}\label{sec:classes-dynamic}

The 1992 manual lets a running program change its classes: \code{add-slot}
and \code{remove-slot} add and remove slots, returning slot descriptors;
\code{(set! (direct-superclasses \meta{class}) \meta{list})} gives a class
new superclasses; and \code{(make <class> superclasses: \meta{classes}
debug-name: \meta{name})} creates a new class. Existing instances are
updated to the new definition before they are next used.

All of this works in the Common Lisp implementation. A program for the
compiler to C is static: its classes and their slots are fixed when it is
compiled. There \code{add-slot} and \code{remove-slot} are errors when
compiling, \code{make} of \code{<class>} is an error when run, and
\code{define-slot} on a class, which the compiler handles as part of the
class's definition, is the way to add a slot (Section~\ref{sec:classes-define-slot}).

%! lisp-only
\begin{dylan}
(define-class <pt> (<object>) (px init-value: 0))
(define p (make <pt>))
(define-generic-function weight (x))
(define-generic-function (setter weight) (x v))
(define sd (add-slot <pt> getter: weight setter: (setter weight)
                          init-value: 3))
(weight p)                     ; => 3
(remove-slot <pt> sd)
(weight p)                     ; => error: no applicable method
(define c (make <class> superclasses: (list <object>) debug-name: '<dyn>))
(instance? (make c) c)         ; => #t
\end{dylan}

\section{The built-in classes}\label{sec:classes-built-in}

The classes below are those of the 1992 manual, as the compiler to C
arranges them. Each class is shown under its first direct superclass; a
class with more than one is marked with the others. Where the Common Lisp
implementation differs, a note says so.

%! skip
\begin{dylan}
<object>
  <boolean>                          (compiler to C only; internal)
  <number>
    <complex>
      <real>
        <rational>
          <integer>
          <ratio>
        <float>
          <single-float>
          <double-float>
          <extended-float>
      <rectangular-complex>          (Common Lisp only)
  <collection>
    <explicit-key-collection>
      <mutable-explicit-key-collection>   (also <mutable-collection>)
        <array>
          <vector>                   (also <mutable-sequence>)
            <simple-object-vector>
            <stretchy-vector>
            <byte-string>            (also <string>)
            <unicode-string>         (also <string>)
        <table>
    <mutable-collection>
      <mutable-sequence>             (also <sequence>)
        <string>
        <list>
          <pair>
          <empty-list>
        <deque>
    <sequence>
      <range>
  <function>
    <generic-function>
    <method>
  <class>
  <singleton>
  <slot-descriptor>
  <symbol>
  <keyword>
  <character>
  <condition>
    <serious-condition>
      <error>
        <simple-error>
        <type-error>
        <thread-error>               (compiler to C only)
    <warning>
      <simple-warning>
    <restart>
      <simple-restart>
      <abort>
\end{dylan}

\noindent Both implementations add the stream class \code{<stream>},
a direct subclass of \code{<object>} with the subclass
\code{<string-stream>} (Chapter~\ref{ch:printing}). The compiler to C
has further classes for its threads (Chapter~\ref{ch:extensions}), also
direct subclasses of \code{<object>}: \code{<thread>}, \code{<lock>},
\code{<condition-variable>} and \code{<atomic-integer>}. \code{<lock>}
has the subclass \code{<recursive-lock>}.

Numbers are described in Chapter~\ref{ch:numbers}, collections in
Chapter~\ref{ch:collections} and conditions in Chapter~\ref{ch:conditions}.

\subsection{Which classes can be subclassed}

A program may define subclasses of \code{<object>} and of the abstract
classes meant to be extended: \code{<number>}, the collection classes
\code{<collection>}, \code{<sequence>}, \code{<mutable-sequence>},
\code{<explicit-key-collection>} and the others above them, the
instantiable collection classes \code{<array>}, \code{<vector>},
\code{<string>}, \code{<table>}, \code{<deque>}, \code{<stretchy-vector>}
and \code{<range>}, and \code{<function>}, \code{<condition>},
\code{<error>} and the other condition classes. These classes are
\emph{sealed}, and subclassing them is an error in both implementations:
the number classes below \code{<number>}, from \code{<complex>} down;
the list classes \code{<list>}, \code{<pair>} and \code{<empty-list>}; and
the concrete vector and string classes \code{<simple-object-vector>},
\code{<byte-string>} and \code{<unicode-string>}.

The compiler to C also refuses, when compiling, to use as a superclass the
classes whose instances its runtime makes itself: \code{<method>},
\code{<generic-function>}, \code{<class>}, \code{<singleton>},
\code{<symbol>}, \code{<keyword>} and \code{<character>}. The Common Lisp
implementation accepts these.
